\documentclass[11pt,a4paper]{article}
% lesson-preamble.tex -- Shared preamble for all Phase 00 lessons
% Usage: % lesson-preamble.tex -- Shared preamble for all Phase 00 lessons
% Usage: % lesson-preamble.tex -- Shared preamble for all Phase 00 lessons
% Usage: \input{../lesson-preamble} (from subdomain directories)

% ---- Encoding and fonts ----
\usepackage[utf8]{inputenc}
\usepackage[T1]{fontenc}

% ---- Mathematics ----
\usepackage{amsmath,amssymb,amsthm}
\usepackage{mathtools}

% ---- Page geometry ----
\usepackage[margin=1in,headheight=14pt]{geometry}

% ---- Hyperlinks ----
\usepackage[colorlinks=true,linkcolor=red,citecolor=blue,urlcolor=blue]{hyperref}

% ---- Lists ----
\usepackage{enumitem}

% ---- Colors and boxes ----
\usepackage{xcolor}
\usepackage[theorems,skins,breakable]{tcolorbox}

% ---- Header/Footer ----
\usepackage{fancyhdr}
\renewcommand{\headrulewidth}{0.4pt}

% ---- Tables ----
\usepackage{booktabs}

% ---- Bibliography ----
\usepackage{natbib}
\bibliographystyle{plainnat}

% --------------------------------------------------------------------------
% Theorem environments
% --------------------------------------------------------------------------
\theoremstyle{plain}
\newtheorem{theorem}{Theorem}[section]
\newtheorem{lemma}[theorem]{Lemma}
\newtheorem{proposition}[theorem]{Proposition}
\newtheorem{corollary}[theorem]{Corollary}

\theoremstyle{definition}
\newtheorem{definition}[theorem]{Definition}
\newtheorem{example}[theorem]{Example}
\newtheorem{exercise}[theorem]{Exercise}
\newtheorem{assumption}[theorem]{Assumption}

\theoremstyle{remark}
\newtheorem{remark}[theorem]{Remark}

% --------------------------------------------------------------------------
% Colored boxes
% --------------------------------------------------------------------------
\newtcolorbox{keyresult}[1][]{%
  colback=green!5!white,
  colframe=green!50!black,
  fonttitle=\bfseries,
  title={Key Result},
  breakable,
  #1
}

\newtcolorbox{rlconnection}[1][]{%
  colback=blue!5!white,
  colframe=blue!50!black,
  fonttitle=\bfseries,
  title={RL/ML Connection},
  breakable,
  #1
}

\newtcolorbox{warning}[1][]{%
  colback=red!5!white,
  colframe=red!50!black,
  fonttitle=\bfseries,
  title={Warning},
  breakable,
  #1
}

% --------------------------------------------------------------------------
% Custom commands -- number sets
% --------------------------------------------------------------------------
\newcommand{\R}{\mathbb{R}}
\newcommand{\N}{\mathbb{N}}
\newcommand{\Q}{\mathbb{Q}}
\newcommand{\Z}{\mathbb{Z}}
\newcommand{\C}{\mathbb{C}}
\newcommand{\F}{\mathbb{F}}

% --------------------------------------------------------------------------
% Custom commands -- probability and expectation
% --------------------------------------------------------------------------
\newcommand{\E}{\mathbb{E}}
\newcommand{\Prob}{\mathbb{P}}
\newcommand{\Var}{\operatorname{Var}}
\newcommand{\Cov}{\operatorname{Cov}}
\newcommand{\indicator}{\mathbf{1}}

% --------------------------------------------------------------------------
% Custom commands -- convergence arrows
% --------------------------------------------------------------------------
\newcommand{\cas}{\xrightarrow{\text{a.s.}}}
\newcommand{\cprob}{\xrightarrow{\Prob}}
\newcommand{\clp}[1]{\xrightarrow{L^{#1}}}
\newcommand{\cdist}{\xrightarrow{d}}

% --------------------------------------------------------------------------
% Custom commands -- common operators
% --------------------------------------------------------------------------
\DeclareMathOperator{\dom}{dom}
\DeclareMathOperator{\epi}{epi}
\DeclareMathOperator{\conv}{conv}
\DeclareMathOperator{\relint}{relint}
\DeclareMathOperator{\aff}{aff}
\DeclareMathOperator{\cl}{cl}
\DeclareMathOperator{\interior}{int}
\DeclareMathOperator{\tr}{tr}
\DeclareMathOperator{\diag}{diag}
\DeclareMathOperator{\rank}{rank}
\DeclareMathOperator{\sgn}{sgn}
\DeclareMathOperator{\supp}{supp}
\DeclareMathOperator{\argmin}{arg\,min}
\DeclareMathOperator{\argmax}{arg\,max}
\DeclareMathOperator{\prox}{prox}
\DeclareMathOperator{\dist}{dist}
\DeclareMathOperator{\diam}{diam}
\DeclareMathOperator{\Span}{span}

% --------------------------------------------------------------------------
% Custom commands -- linear algebra operators
% --------------------------------------------------------------------------
\DeclareMathOperator{\nullity}{nullity}
\DeclareMathOperator{\range}{range}
\DeclareMathOperator{\nullspace}{null}
\DeclareMathOperator{\proj}{proj}

% --------------------------------------------------------------------------
% Custom commands -- measure theory
% --------------------------------------------------------------------------
\newcommand{\powerset}{\mathcal{P}}
\newcommand{\Borel}{\mathcal{B}}
 (from subdomain directories)

% ---- Encoding and fonts ----
\usepackage[utf8]{inputenc}
\usepackage[T1]{fontenc}

% ---- Mathematics ----
\usepackage{amsmath,amssymb,amsthm}
\usepackage{mathtools}

% ---- Page geometry ----
\usepackage[margin=1in,headheight=14pt]{geometry}

% ---- Hyperlinks ----
\usepackage[colorlinks=true,linkcolor=red,citecolor=blue,urlcolor=blue]{hyperref}

% ---- Lists ----
\usepackage{enumitem}

% ---- Colors and boxes ----
\usepackage{xcolor}
\usepackage[theorems,skins,breakable]{tcolorbox}

% ---- Header/Footer ----
\usepackage{fancyhdr}
\renewcommand{\headrulewidth}{0.4pt}

% ---- Tables ----
\usepackage{booktabs}

% ---- Bibliography ----
\usepackage{natbib}
\bibliographystyle{plainnat}

% --------------------------------------------------------------------------
% Theorem environments
% --------------------------------------------------------------------------
\theoremstyle{plain}
\newtheorem{theorem}{Theorem}[section]
\newtheorem{lemma}[theorem]{Lemma}
\newtheorem{proposition}[theorem]{Proposition}
\newtheorem{corollary}[theorem]{Corollary}

\theoremstyle{definition}
\newtheorem{definition}[theorem]{Definition}
\newtheorem{example}[theorem]{Example}
\newtheorem{exercise}[theorem]{Exercise}
\newtheorem{assumption}[theorem]{Assumption}

\theoremstyle{remark}
\newtheorem{remark}[theorem]{Remark}

% --------------------------------------------------------------------------
% Colored boxes
% --------------------------------------------------------------------------
\newtcolorbox{keyresult}[1][]{%
  colback=green!5!white,
  colframe=green!50!black,
  fonttitle=\bfseries,
  title={Key Result},
  breakable,
  #1
}

\newtcolorbox{rlconnection}[1][]{%
  colback=blue!5!white,
  colframe=blue!50!black,
  fonttitle=\bfseries,
  title={RL/ML Connection},
  breakable,
  #1
}

\newtcolorbox{warning}[1][]{%
  colback=red!5!white,
  colframe=red!50!black,
  fonttitle=\bfseries,
  title={Warning},
  breakable,
  #1
}

% --------------------------------------------------------------------------
% Custom commands -- number sets
% --------------------------------------------------------------------------
\newcommand{\R}{\mathbb{R}}
\newcommand{\N}{\mathbb{N}}
\newcommand{\Q}{\mathbb{Q}}
\newcommand{\Z}{\mathbb{Z}}
\newcommand{\C}{\mathbb{C}}
\newcommand{\F}{\mathbb{F}}

% --------------------------------------------------------------------------
% Custom commands -- probability and expectation
% --------------------------------------------------------------------------
\newcommand{\E}{\mathbb{E}}
\newcommand{\Prob}{\mathbb{P}}
\newcommand{\Var}{\operatorname{Var}}
\newcommand{\Cov}{\operatorname{Cov}}
\newcommand{\indicator}{\mathbf{1}}

% --------------------------------------------------------------------------
% Custom commands -- convergence arrows
% --------------------------------------------------------------------------
\newcommand{\cas}{\xrightarrow{\text{a.s.}}}
\newcommand{\cprob}{\xrightarrow{\Prob}}
\newcommand{\clp}[1]{\xrightarrow{L^{#1}}}
\newcommand{\cdist}{\xrightarrow{d}}

% --------------------------------------------------------------------------
% Custom commands -- common operators
% --------------------------------------------------------------------------
\DeclareMathOperator{\dom}{dom}
\DeclareMathOperator{\epi}{epi}
\DeclareMathOperator{\conv}{conv}
\DeclareMathOperator{\relint}{relint}
\DeclareMathOperator{\aff}{aff}
\DeclareMathOperator{\cl}{cl}
\DeclareMathOperator{\interior}{int}
\DeclareMathOperator{\tr}{tr}
\DeclareMathOperator{\diag}{diag}
\DeclareMathOperator{\rank}{rank}
\DeclareMathOperator{\sgn}{sgn}
\DeclareMathOperator{\supp}{supp}
\DeclareMathOperator{\argmin}{arg\,min}
\DeclareMathOperator{\argmax}{arg\,max}
\DeclareMathOperator{\prox}{prox}
\DeclareMathOperator{\dist}{dist}
\DeclareMathOperator{\diam}{diam}
\DeclareMathOperator{\Span}{span}

% --------------------------------------------------------------------------
% Custom commands -- linear algebra operators
% --------------------------------------------------------------------------
\DeclareMathOperator{\nullity}{nullity}
\DeclareMathOperator{\range}{range}
\DeclareMathOperator{\nullspace}{null}
\DeclareMathOperator{\proj}{proj}

% --------------------------------------------------------------------------
% Custom commands -- measure theory
% --------------------------------------------------------------------------
\newcommand{\powerset}{\mathcal{P}}
\newcommand{\Borel}{\mathcal{B}}
 (from subdomain directories)

% ---- Encoding and fonts ----
\usepackage[utf8]{inputenc}
\usepackage[T1]{fontenc}

% ---- Mathematics ----
\usepackage{amsmath,amssymb,amsthm}
\usepackage{mathtools}

% ---- Page geometry ----
\usepackage[margin=1in,headheight=14pt]{geometry}

% ---- Hyperlinks ----
\usepackage[colorlinks=true,linkcolor=red,citecolor=blue,urlcolor=blue]{hyperref}

% ---- Lists ----
\usepackage{enumitem}

% ---- Colors and boxes ----
\usepackage{xcolor}
\usepackage[theorems,skins,breakable]{tcolorbox}

% ---- Header/Footer ----
\usepackage{fancyhdr}
\renewcommand{\headrulewidth}{0.4pt}

% ---- Tables ----
\usepackage{booktabs}

% ---- Bibliography ----
\usepackage{natbib}
\bibliographystyle{plainnat}

% --------------------------------------------------------------------------
% Theorem environments
% --------------------------------------------------------------------------
\theoremstyle{plain}
\newtheorem{theorem}{Theorem}[section]
\newtheorem{lemma}[theorem]{Lemma}
\newtheorem{proposition}[theorem]{Proposition}
\newtheorem{corollary}[theorem]{Corollary}

\theoremstyle{definition}
\newtheorem{definition}[theorem]{Definition}
\newtheorem{example}[theorem]{Example}
\newtheorem{exercise}[theorem]{Exercise}
\newtheorem{assumption}[theorem]{Assumption}

\theoremstyle{remark}
\newtheorem{remark}[theorem]{Remark}

% --------------------------------------------------------------------------
% Colored boxes
% --------------------------------------------------------------------------
\newtcolorbox{keyresult}[1][]{%
  colback=green!5!white,
  colframe=green!50!black,
  fonttitle=\bfseries,
  title={Key Result},
  breakable,
  #1
}

\newtcolorbox{rlconnection}[1][]{%
  colback=blue!5!white,
  colframe=blue!50!black,
  fonttitle=\bfseries,
  title={RL/ML Connection},
  breakable,
  #1
}

\newtcolorbox{warning}[1][]{%
  colback=red!5!white,
  colframe=red!50!black,
  fonttitle=\bfseries,
  title={Warning},
  breakable,
  #1
}

% --------------------------------------------------------------------------
% Custom commands -- number sets
% --------------------------------------------------------------------------
\newcommand{\R}{\mathbb{R}}
\newcommand{\N}{\mathbb{N}}
\newcommand{\Q}{\mathbb{Q}}
\newcommand{\Z}{\mathbb{Z}}
\newcommand{\C}{\mathbb{C}}
\newcommand{\F}{\mathbb{F}}

% --------------------------------------------------------------------------
% Custom commands -- probability and expectation
% --------------------------------------------------------------------------
\newcommand{\E}{\mathbb{E}}
\newcommand{\Prob}{\mathbb{P}}
\newcommand{\Var}{\operatorname{Var}}
\newcommand{\Cov}{\operatorname{Cov}}
\newcommand{\indicator}{\mathbf{1}}

% --------------------------------------------------------------------------
% Custom commands -- convergence arrows
% --------------------------------------------------------------------------
\newcommand{\cas}{\xrightarrow{\text{a.s.}}}
\newcommand{\cprob}{\xrightarrow{\Prob}}
\newcommand{\clp}[1]{\xrightarrow{L^{#1}}}
\newcommand{\cdist}{\xrightarrow{d}}

% --------------------------------------------------------------------------
% Custom commands -- common operators
% --------------------------------------------------------------------------
\DeclareMathOperator{\dom}{dom}
\DeclareMathOperator{\epi}{epi}
\DeclareMathOperator{\conv}{conv}
\DeclareMathOperator{\relint}{relint}
\DeclareMathOperator{\aff}{aff}
\DeclareMathOperator{\cl}{cl}
\DeclareMathOperator{\interior}{int}
\DeclareMathOperator{\tr}{tr}
\DeclareMathOperator{\diag}{diag}
\DeclareMathOperator{\rank}{rank}
\DeclareMathOperator{\sgn}{sgn}
\DeclareMathOperator{\supp}{supp}
\DeclareMathOperator{\argmin}{arg\,min}
\DeclareMathOperator{\argmax}{arg\,max}
\DeclareMathOperator{\prox}{prox}
\DeclareMathOperator{\dist}{dist}
\DeclareMathOperator{\diam}{diam}
\DeclareMathOperator{\Span}{span}

% --------------------------------------------------------------------------
% Custom commands -- linear algebra operators
% --------------------------------------------------------------------------
\DeclareMathOperator{\nullity}{nullity}
\DeclareMathOperator{\range}{range}
\DeclareMathOperator{\nullspace}{null}
\DeclareMathOperator{\proj}{proj}

% --------------------------------------------------------------------------
% Custom commands -- measure theory
% --------------------------------------------------------------------------
\newcommand{\powerset}{\mathcal{P}}
\newcommand{\Borel}{\mathcal{B}}

\pagestyle{fancy}
\fancyhf{}
\fancyhead[L]{\textsc{Lesson 8: SVD, Norms, and Matrix Calculus}}
\fancyhead[R]{\thepage}
\title{Lesson 8: SVD, Norms, and Matrix Calculus}
\author{AI/ML Mathematics}
\date{}
\begin{document}
\maketitle
\tableofcontents
\newpage

%% ============================================================
%% SECTION 1: MOTIVATION AND OVERVIEW
%% ============================================================
\section{Motivation and Overview}
\label{sec:motivation}

How does a neural network learn?  At each training step, the algorithm
computes the gradient of a scalar loss function with respect to a
matrix of weights, updates the weights by a small step in the negative
gradient direction, and repeats.  Unpacking this single sentence
reveals three distinct mathematical needs.  First, the weight matrix
$W \in \R^{m \times n}$ transforms input features, and understanding
its action -- which directions it amplifies, which it suppresses, and
what its effective rank is -- requires the \emph{singular value
decomposition} (SVD).  Second, bounding the sensitivity of the network
output to perturbations in $W$ requires \emph{matrix norms}, and the
relationship between the operator norm and the Frobenius norm governs
both the stability of the forward pass and the explosion or vanishing
of gradients during backpropagation.  Third, computing the gradient
itself -- a derivative of a scalar with respect to a matrix -- belongs
to \emph{matrix calculus}, the systematic extension of multivariable
calculus to matrix-valued arguments.

This capstone lesson of the Linear Algebra subdomain assembles
these three pillars: SVD, matrix norms, and matrix calculus.  Along
the way, we develop quadratic forms (the local shape of any smooth loss
near a critical point), the matrix exponential (the bridge to
continuous-time dynamics), and the derivative formulas that underpin
every gradient computation in deep learning and reinforcement learning.

\begin{rlconnection}[title={Where SVD, Norms, and Matrix Calculus Appear in RL/ML}]
\begin{itemize}[nosep]
  \item \textbf{PCA and dimensionality reduction.}
        Principal component analysis computes the top singular vectors
        of the centered data matrix; the Eckart--Young theorem guarantees
        this is the best low-rank approximation in both operator and
        Frobenius norm.
  \item \textbf{Backpropagation and gradient computation.}
        Every step of backpropagation applies the chain rule for
        matrix-valued functions; the formulas $\nabla_x(x^\top A x)
        = (A + A^\top)x$ and $\partial \tr(AX)/\partial X = A^\top$
        appear in nearly every derivation.
  \item \textbf{Condition number and optimization convergence.}
        The condition number $\kappa(A) = \sigma_1/\sigma_n$ controls
        the convergence rate of gradient descent; ill-conditioned
        Hessians cause slow training, motivating preconditioning and
        adaptive methods like Adam.
  \item \textbf{Policy gradient and natural gradient.}
        The natural policy gradient requires differentiating the
        KL-divergence with respect to parameters, which involves
        quadratic forms in the Fisher information matrix and the
        matrix calculus chain rule.
\end{itemize}
\end{rlconnection}

\paragraph{Prerequisites.}
Vector spaces, bases, dimension, linear maps, rank--nullity theorem
(Lesson~5).  Determinants, eigenvalues, characteristic polynomial,
diagonalizability criterion, Cayley--Hamilton theorem (Lesson~6).
Inner products, orthonormal bases, Gram--Schmidt, orthogonal
projection, spectral theorem for real symmetric matrices, positive
definiteness, Rayleigh quotient (Lesson~7).

\paragraph{Learning objectives.}
After completing this lesson, the student will be able to:
\begin{itemize}[nosep]
  \item \textbf{Classify} quadratic forms by the eigenvalues of the
        associated symmetric matrix and connect them to second-order
        optimality conditions.
  \item \textbf{Prove} the SVD theorem from the spectral decomposition
        of $A^\top A$ and derive the Eckart--Young low-rank approximation
        theorem.
  \item \textbf{Compute} the operator norm and Frobenius norm of a matrix
        in terms of singular values and prove their key relationships.
  \item \textbf{Derive} the matrix exponential and its fundamental
        properties, including the connection to linear ODEs.
  \item \textbf{Apply} the rules of matrix calculus -- gradient,
        Jacobian, Hessian, trace derivatives, and the chain rule --
        to compute derivatives arising in optimization and machine
        learning.
  \item \textbf{Construct} complete gradient derivations for quadratic
        objectives and trace-based loss functions using the tools of
        this lesson.
\end{itemize}

%% ============================================================
%% SECTION 2: REFERENCES
%% ============================================================
\section{References}
\label{sec:references}

\begin{itemize}
  \item \citet[Chapter 7]{axler2015} -- Singular value decomposition
        from the operator-theoretic perspective, polar decomposition.
  \item \citet[Sections 3.1--3.3, 7.4]{horn2013} -- Comprehensive
        treatment of SVD, matrix norms, condition numbers, and the
        Eckart--Young theorem.
  \item \citet[Chapters 6--7]{strang2016} -- Computational perspective
        on SVD, least squares, pseudoinverse, and matrix norms.
  \item \citet[Chapters 1--6]{magnus2019} -- Rigorous development of
        matrix differential calculus with applications.
  \item \citet{petersen2012} -- Extensive catalog of matrix derivative
        identities organized by matrix structure.
  \item \citet[Chapter 6]{lax2007} -- Matrix exponential, linear ODEs,
        and the connection to semigroup theory.
\end{itemize}

%% ============================================================
%% SECTION 3: QUADRATIC FORMS
%% ============================================================
\section{Quadratic Forms}
\label{sec:quadratic_forms}

\subsection{Definition and Canonical Representation}

\begin{definition}[Quadratic Form]
\label{def:quadratic_form}
Let $A \in \R^{n \times n}$ be a symmetric matrix.  The
\emph{quadratic form} associated with $A$ is the function
$Q_A\colon \R^n \to \R$ defined by
\[
  Q_A(x) = x^\top A x = \sum_{i=1}^n \sum_{j=1}^n a_{ij} x_i x_j.
\]
\end{definition}

\begin{remark}
\label{rem:symmetrize}
Any matrix $B \in \R^{n \times n}$ (not necessarily symmetric) gives
a quadratic form $x^\top B x = x^\top \bigl(\frac{B + B^\top}{2}\bigr) x$,
since $x^\top B x$ is a scalar and therefore equals its own transpose
$x^\top B^\top x$.  Hence every quadratic form can be represented by a
unique symmetric matrix $A = (B + B^\top)/2$.
\end{remark}

\begin{definition}[Classification of Quadratic Forms]
\label{def:classification_qf}
A quadratic form $Q_A$ (equivalently, its symmetric matrix $A$) is:
\begin{itemize}[nosep]
  \item \emph{positive definite} if $Q_A(x) > 0$ for all $x \neq 0$,
  \item \emph{positive semidefinite} if $Q_A(x) \geq 0$ for all $x$,
  \item \emph{negative definite} if $Q_A(x) < 0$ for all $x \neq 0$,
  \item \emph{negative semidefinite} if $Q_A(x) \leq 0$ for all $x$,
  \item \emph{indefinite} if $Q_A$ takes both positive and negative values.
\end{itemize}
\end{definition}

\begin{proposition}[Eigenvalue Classification]
\label{prop:eigenvalue_classification}
Let $A \in \R^{n \times n}$ be symmetric with eigenvalues
$\lambda_1, \ldots, \lambda_n$.  Then:
\begin{enumerate}[nosep]
  \item $A \succ 0$ if and only if $\lambda_i > 0$ for all $i$.
  \item $A \succeq 0$ if and only if $\lambda_i \geq 0$ for all $i$.
  \item $A$ is indefinite if and only if $A$ has both positive and
        negative eigenvalues.
\end{enumerate}
\end{proposition}

\begin{proof}
By the spectral theorem (Lesson~7, Theorem~5.3), $A = Q \Lambda Q^\top$
where $Q$ is orthogonal and $\Lambda = \diag(\lambda_1, \ldots, \lambda_n)$.
Setting $y = Q^\top x$ (a bijection since $Q$ is invertible):
\[
  Q_A(x) = x^\top Q \Lambda Q^\top x = y^\top \Lambda y
  = \sum_{i=1}^n \lambda_i y_i^2.
\]
If all $\lambda_i > 0$, then $\sum_i \lambda_i y_i^2 > 0$ for $y \neq 0$.
Conversely, if $\lambda_k \leq 0$ for some $k$, then $x = q_k$ gives
$Q_A(q_k) = \lambda_k \leq 0$.  This proves (1).  Parts (2) and (3)
follow by the same argument.
\end{proof}

\begin{example}[Classifying a Quadratic Form]
\label{ex:classify_qf}
Let $A = \begin{pmatrix} 2 & -1 \\ -1 & 2 \end{pmatrix}$.  The
eigenvalues are $\lambda = 2 \pm 1$, i.e., $\lambda_1 = 1$ and
$\lambda_2 = 3$.  Since both are positive, $A \succ 0$ and the
quadratic form $Q_A(x) = 2x_1^2 - 2x_1 x_2 + 2x_2^2$ is positive
definite.
\end{example}

\subsection{Quadratic Approximation and Second-Order Conditions}

\begin{theorem}[Second-Order Taylor Expansion]
\label{thm:taylor_second_order}
Let $f\colon \R^n \to \R$ be twice continuously differentiable.  At a
point $x_0 \in \R^n$, for any $h \in \R^n$:
\[
  f(x_0 + h) = f(x_0) + \nabla f(x_0)^\top h
  + \frac{1}{2} h^\top H_f(x_0)\, h + o(\|h\|^2),
\]
where $H_f(x_0) \in \R^{n \times n}$ is the Hessian matrix with entries
$[H_f(x_0)]_{ij} = \frac{\partial^2 f}{\partial x_i \partial x_j}(x_0)$.
\end{theorem}

\begin{proof}
Apply the one-dimensional Taylor expansion to $g(t) = f(x_0 + th)$
at $t = 0$: $g(1) = g(0) + g'(0) + \frac{1}{2}g''(0) + o(1)$ as
$\|h\| \to 0$.  By the chain rule, $g'(0) = \nabla f(x_0)^\top h$
and $g''(0) = h^\top H_f(x_0)\, h$.
\end{proof}

\begin{corollary}[Second-Order Optimality Conditions]
\label{cor:second_order_opt}
Let $f\colon \R^n \to \R$ be twice continuously differentiable and let
$x_0$ be a critical point ($\nabla f(x_0) = 0$).  Then:
\begin{enumerate}[nosep]
  \item If $H_f(x_0) \succ 0$, then $x_0$ is a strict local minimum.
  \item If $H_f(x_0) \prec 0$, then $x_0$ is a strict local maximum.
  \item If $H_f(x_0)$ is indefinite, then $x_0$ is a saddle point.
\end{enumerate}
\end{corollary}

\begin{proof}
At a critical point, $f(x_0 + h) - f(x_0) = \frac{1}{2}h^\top H_f(x_0)\,h
+ o(\|h\|^2)$.  If $H_f(x_0) \succ 0$, the quadratic term dominates
for small $\|h\|$ and is positive for $h \neq 0$, giving a strict local
minimum.  Similarly for the other cases.
\end{proof}

\begin{example}[Quadratic Approximation Near a Minimum]
\label{ex:quad_approx}
Let $f(x_1, x_2) = x_1^4 + x_1^2 + x_2^2$.  At $x_0 = (0,0)$:
$\nabla f(0) = (0, 0)^\top$ and
$H_f(0) = \begin{pmatrix} 2 & 0 \\ 0 & 2 \end{pmatrix} = 2I$.
Since $H_f(0) \succ 0$, the origin is a strict local minimum.
The quadratic approximation is $f(h) \approx h_1^2 + h_2^2 = \|h\|^2$.
\end{example}

\begin{rlconnection}[title={Newton's Method and the Natural Policy Gradient}]
At a critical point, Newton's method for minimization updates
$x \leftarrow x - H_f(x)^{-1}\nabla f(x)$, effectively using the
quadratic form $h^\top H_f(x)\,h$ to determine both the direction and
step size.  In policy optimization, the natural policy gradient replaces
the Euclidean Hessian with the Fisher information matrix
$F(\theta) \succeq 0$, yielding the update
$\theta \leftarrow \theta + \alpha F(\theta)^{-1}\nabla J(\theta)$.
The Fisher matrix defines a quadratic form $h^\top F(\theta)\, h$
that measures the local curvature of the KL-divergence between nearby
policies.  The classification of $F(\theta)$ as positive semidefinite
guarantees that this curvature measure is well-defined.
\end{rlconnection}

With quadratic forms providing the local geometry of smooth functions,
we now turn to a factorization that reveals the global geometry of
any linear map: the singular value decomposition.

%% ============================================================
%% SECTION 4: SINGULAR VALUE DECOMPOSITION
%% ============================================================
\section{Singular Value Decomposition}
\label{sec:svd}

\subsection{Statement and Proof}

\begin{definition}[Singular Values]
\label{def:singular_values}
Let $A \in \R^{m \times n}$.  The \emph{singular values} of $A$ are
$\sigma_1 \geq \sigma_2 \geq \cdots \geq \sigma_{\min(m,n)} \geq 0$,
defined as the square roots of the eigenvalues of $A^\top A \in \R^{n \times n}$
(or equivalently of the nonzero eigenvalues of $AA^\top \in \R^{m \times m}$).
\end{definition}

\begin{remark}
\label{rem:ata_psd}
The matrix $A^\top A$ is symmetric and positive semidefinite, since
$x^\top(A^\top A)x = \|Ax\|^2 \geq 0$ for all $x$.  Therefore all
eigenvalues of $A^\top A$ are nonneg, and the singular values are
well-defined real numbers.
\end{remark}

\begin{keyresult}[title={The Singular Value Decomposition}]
\begin{theorem}[SVD Theorem]
\label{thm:svd}
Let $A \in \R^{m \times n}$ with $\rank(A) = r$.  There exist an
orthogonal matrix $U \in \R^{m \times m}$, an orthogonal matrix
$V \in \R^{n \times n}$, and a matrix $\Sigma \in \R^{m \times n}$
with $\Sigma_{ii} = \sigma_i$ for $i = 1, \ldots, \min(m,n)$ and all
other entries zero, such that
\[
  A = U \Sigma V^\top.
\]
The singular values satisfy
$\sigma_1 \geq \cdots \geq \sigma_r > 0 = \sigma_{r+1} = \cdots = \sigma_{\min(m,n)}$.
The columns $u_1, \ldots, u_m$ of $U$ are called the
\emph{left singular vectors}, and the columns $v_1, \ldots, v_n$
of $V$ are called the \emph{right singular vectors}.
\end{theorem}
\end{keyresult}

\begin{proof}
Since $A^\top A \in \R^{n \times n}$ is symmetric positive semidefinite,
the spectral theorem (Lesson~7, Theorem~5.3) provides an orthogonal
matrix $V = (v_1 \mid \cdots \mid v_n)$ and eigenvalues
$\mu_1 \geq \cdots \geq \mu_n \geq 0$ such that
\[
  A^\top A = V \diag(\mu_1, \ldots, \mu_n) V^\top.
\]
Set $\sigma_i = \sqrt{\mu_i}$ for each $i$.  Since
$\rank(A^\top A) = \rank(A) = r$ (the rank of $A^\top A$ equals the
rank of $A$, as $\nullspace(A^\top A) = \nullspace(A)$), we have
$\sigma_1 \geq \cdots \geq \sigma_r > 0$ and
$\sigma_{r+1} = \cdots = \sigma_n = 0$.

\textbf{Step 1: Construct the left singular vectors for nonzero
singular values.}  For $i = 1, \ldots, r$, define
\[
  u_i = \frac{1}{\sigma_i} A v_i.
\]
We verify orthonormality: for $1 \leq i, j \leq r$,
\[
  u_i^\top u_j = \frac{1}{\sigma_i \sigma_j} v_i^\top A^\top A v_j
  = \frac{1}{\sigma_i \sigma_j} v_i^\top (\sigma_j^2 v_j)
  = \frac{\sigma_j}{\sigma_i} v_i^\top v_j
  = \frac{\sigma_j}{\sigma_i} \delta_{ij} = \delta_{ij},
\]
where we used $A^\top A v_j = \sigma_j^2 v_j$ and the orthonormality
of the $v_j$.

\textbf{Step 2: Extend to an orthonormal basis of $\R^m$.}
The vectors $u_1, \ldots, u_r$ are orthonormal in $\R^m$.  If $r < m$,
extend them to an orthonormal basis $u_1, \ldots, u_m$ of $\R^m$
(possible by the Gram--Schmidt process applied to any completion).
Set $U = (u_1 \mid \cdots \mid u_m)$.

\textbf{Step 3: Verify $A = U\Sigma V^\top$.}  It suffices to check
$AV = U\Sigma$.  For $i = 1, \ldots, r$:
$Av_i = \sigma_i u_i = (U\Sigma)_{\cdot, i}$.
For $i = r+1, \ldots, n$: $\|Av_i\|^2 = v_i^\top A^\top A v_i
= \sigma_i^2 \|v_i\|^2 = 0$, so $Av_i = 0 = (U\Sigma)_{\cdot, i}$
(since $\sigma_i = 0$).  Thus $AV = U\Sigma$, giving
$A = U\Sigma V^\top$.
\end{proof}

\subsection{Geometric Interpretation}

The SVD reveals that every linear map $A\colon \R^n \to \R^m$ acts
by three steps: (1) rotate/reflect by $V^\top$ in the domain, (2) scale
by $\sigma_i$ along each coordinate axis (and embed or project if
$m \neq n$), (3) rotate/reflect by $U$ in the codomain.  In particular,
the image of the unit sphere $\{x : \|x\| = 1\}$ under $A$ is an
ellipsoid in $\R^m$ with semi-axes $\sigma_1 u_1, \ldots, \sigma_r u_r$.

\begin{example}[SVD of a $2 \times 2$ Matrix]
\label{ex:svd_2x2}
Let $A = \begin{pmatrix} 3 & 0 \\ 0 & 2 \end{pmatrix}$.
Then $A^\top A = \begin{pmatrix} 9 & 0 \\ 0 & 4 \end{pmatrix}$, with
eigenvalues $\mu_1 = 9$, $\mu_2 = 4$ and eigenvectors $v_1 = e_1$,
$v_2 = e_2$.  So $\sigma_1 = 3$, $\sigma_2 = 2$, $u_1 = Av_1/3 = e_1$,
$u_2 = Av_2/2 = e_2$.  The SVD is $A = I \cdot \diag(3,2) \cdot I^\top$,
which is the matrix itself -- diagonal matrices are already in SVD form.
\end{example}

\begin{example}[SVD of a Non-Square Matrix]
\label{ex:svd_nonsquare}
Let $A = \begin{pmatrix} 1 & 1 \\ 0 & 0 \end{pmatrix}$.
Then $A^\top A = \begin{pmatrix} 1 & 1 \\ 1 & 1 \end{pmatrix}$,
with eigenvalues $\mu_1 = 2$, $\mu_2 = 0$ and eigenvectors
$v_1 = \frac{1}{\sqrt{2}}(1,1)^\top$,
$v_2 = \frac{1}{\sqrt{2}}(1,-1)^\top$.
So $\sigma_1 = \sqrt{2}$, $\sigma_2 = 0$, and
$u_1 = Av_1/\sqrt{2} = (1, 0)^\top = e_1$.
Extending: $u_2 = e_2$.  The SVD is
$A = \begin{pmatrix} 1 & 0 \\ 0 & 1 \end{pmatrix}
\begin{pmatrix} \sqrt{2} & 0 \\ 0 & 0 \end{pmatrix}
\frac{1}{\sqrt{2}}\begin{pmatrix} 1 & 1 \\ 1 & -1 \end{pmatrix}^\top$.
The rank is $r = 1$, confirmed by $\sigma_1 > 0 = \sigma_2$.
\end{example}

\begin{proposition}[Rank from SVD]
\label{prop:rank_svd}
$\rank(A)$ equals the number of nonzero singular values of $A$.
\end{proposition}

\begin{proof}
From the SVD, $\rank(A) = \rank(U\Sigma V^\top) = \rank(\Sigma)$
(since $U$ and $V$ are invertible), and $\rank(\Sigma)$ is the number
of nonzero diagonal entries.
\end{proof}

\subsection{The Eckart--Young Theorem}

\begin{keyresult}[title={Best Low-Rank Approximation}]
\begin{theorem}[Eckart--Young]
\label{thm:eckart_young}
Let $A \in \R^{m \times n}$ have SVD $A = \sum_{i=1}^r \sigma_i u_i v_i^\top$.
For any $k \leq r$, define $A_k = \sum_{i=1}^k \sigma_i u_i v_i^\top$.
Then $A_k$ is the best rank-$k$ approximation to $A$ in both the
operator norm and the Frobenius norm:
\begin{align}
  \min_{\rank(B) \leq k} \|A - B\|_{\mathrm{op}} &= \|A - A_k\|_{\mathrm{op}} = \sigma_{k+1},
  \label{eq:eckart_young_op} \\
  \min_{\rank(B) \leq k} \|A - B\|_{F} &= \|A - A_k\|_{F}
  = \Bigl(\sum_{i=k+1}^r \sigma_i^2\Bigr)^{1/2}.
  \label{eq:eckart_young_frob}
\end{align}
(The norms $\|\cdot\|_{\mathrm{op}}$ and $\|\cdot\|_F$ are defined
formally in Section~\ref{sec:norms}.)
\end{theorem}
\end{keyresult}

\begin{proof}
We prove the operator norm case; the Frobenius norm case follows by a
similar argument.

\textbf{Achievability:} $A - A_k = \sum_{i=k+1}^r \sigma_i u_i v_i^\top$,
which has singular values $\sigma_{k+1} \geq \cdots \geq \sigma_r$.
Hence $\|A - A_k\|_{\mathrm{op}} = \sigma_{k+1}$ (by
Proposition~\ref{prop:op_norm_sigma} below).

\textbf{Lower bound:} Let $B \in \R^{m \times n}$ with
$\rank(B) \leq k$.  Then $\nullspace(B)$ has dimension
$\geq n - k$.  Let $S = \Span(v_1, \ldots, v_{k+1})$, which has
dimension $k + 1$.  Since $\dim(\nullspace(B)) + \dim(S) \geq
(n - k) + (k + 1) = n + 1 > n$, there exists a unit vector
$z \in \nullspace(B) \cap S$.

Write $z = \sum_{i=1}^{k+1} c_i v_i$ with $\sum_{i=1}^{k+1} c_i^2 = 1$.
Then $Bz = 0$ and
\[
  \|A - B\|_{\mathrm{op}} \geq \|(A - B)z\| = \|Az\|
  = \Bigl\|\sum_{i=1}^{k+1} c_i \sigma_i u_i\Bigr\|
  = \Bigl(\sum_{i=1}^{k+1} c_i^2 \sigma_i^2\Bigr)^{1/2}
  \geq \sigma_{k+1},
\]
where the last inequality uses $\sigma_i \geq \sigma_{k+1}$ for
$i \leq k+1$ and $\sum c_i^2 = 1$.

For the Frobenius norm, a similar dimension argument applied to each
singular direction establishes $\|A - B\|_F^2 \geq \sum_{i=k+1}^r \sigma_i^2$
for any rank-$k$ matrix $B$, with equality at $B = A_k$.
\end{proof}

\subsection{The Moore--Penrose Pseudoinverse}

\begin{definition}[Moore--Penrose Pseudoinverse]
\label{def:pseudoinverse}
Let $A \in \R^{m \times n}$ have SVD
$A = \sum_{i=1}^r \sigma_i u_i v_i^\top$.  The \emph{Moore--Penrose
pseudoinverse} of $A$ is
\[
  A^+ = \sum_{i=1}^r \frac{1}{\sigma_i} v_i u_i^\top \in \R^{n \times m}.
\]
Equivalently, if $A = U\Sigma V^\top$, then
$A^+ = V\Sigma^+ U^\top$, where $\Sigma^+$ is obtained from $\Sigma$
by transposing and inverting all nonzero entries.
\end{definition}

\begin{theorem}[Penrose Conditions]
\label{thm:penrose}
The pseudoinverse $A^+$ is the unique matrix $X \in \R^{n \times m}$
satisfying:
\begin{enumerate}[nosep]
  \item $AXA = A$,
  \item $XAX = X$,
  \item $(AX)^\top = AX$,
  \item $(XA)^\top = XA$.
\end{enumerate}
\end{theorem}

\begin{proof}
\textbf{$A^+$ satisfies the four conditions.}
Using $A = U\Sigma V^\top$ and $X = V\Sigma^+ U^\top$:
$AX = U\Sigma\Sigma^+ U^\top$ and $XA = V\Sigma^+\Sigma V^\top$.
Both $\Sigma\Sigma^+$ and $\Sigma^+\Sigma$ are diagonal with $1$'s
in the first $r$ positions and $0$'s elsewhere.  Hence $AX$ and $XA$
are symmetric (conditions 3--4), and $AXA = U(\Sigma\Sigma^+\Sigma)V^\top
= U\Sigma V^\top = A$ (condition 1), $XAX = V(\Sigma^+\Sigma\Sigma^+)U^\top
= V\Sigma^+ U^\top = X$ (condition 2).

\textbf{Uniqueness.}
Let $X$ and $Y$ both satisfy conditions 1--4.  Then
\begin{align*}
  X &= XAX = X(AX)^\top = XX^\top A^\top = XX^\top(AYA)^\top
  = XX^\top A^\top Y^\top A^\top \\
  &= (AX)^\top (AY)^\top A^\top = (AX)^\top [(AY)A]^\top
  = (AX)^\top A^\top Y^\top \\
  &= (AAX)^\top Y^\top = (A)^\top Y^\top = (YA)^\top Y^\top \cdot\,\text{?}
\end{align*}
A cleaner route: from condition 1, $AXA = A$ and $AYA = A$.  Then
$AX$ and $AY$ are both orthogonal projections onto $\range(A)$ (symmetric
and idempotent by conditions 1, 3), so $AX = AY$.  Similarly, $XA = YA$
(projections onto $\range(A^\top)$).  Then
$X = XAX = XAY = (XA)Y = (YA)Y = YAY = Y$.
\end{proof}

\begin{theorem}[Pseudoinverse and Least Squares]
\label{thm:pseudo_ls}
For any $A \in \R^{m \times n}$ and $b \in \R^m$, $x^* = A^+ b$ is
the minimum-norm least-squares solution: among all vectors minimizing
$\|Ax - b\|$, the vector $x^*$ has the smallest $\|x\|$.
\end{theorem}

\begin{proof}
The least-squares problem minimizes $\|Ax - b\|^2$, which is equivalent
to solving the normal equations $A^\top A x = A^\top b$.  The solution
set is $\{x : Ax = \proj_{\range(A)} b\} = \{A^+ b + z : z \in \nullspace(A)\}$.

To see that $x^* = A^+ b$ has minimum norm among all solutions,
note that $A^+ b \in \range(A^\top) = \range(V_r)$ where
$V_r = (v_1 \mid \cdots \mid v_r)$, since
$A^+ b = \sum_{i=1}^r \frac{u_i^\top b}{\sigma_i} v_i$.
Any other solution is $x = x^* + z$ with $z \in \nullspace(A)$.
Since $\range(A^\top) \perp \nullspace(A)$
(the fundamental theorem of linear algebra), the Pythagorean theorem
gives $\|x\|^2 = \|x^*\|^2 + \|z\|^2 \geq \|x^*\|^2$.
\end{proof}

\begin{example}[Pseudoinverse of a Rank-1 Matrix]
\label{ex:pseudo_rank1}
Let $A = \begin{pmatrix} 1 \\ 2 \end{pmatrix} \in \R^{2 \times 1}$.
The SVD is $A = u_1 \sigma_1 v_1^\top$ with $\sigma_1 = \sqrt{5}$,
$u_1 = \frac{1}{\sqrt{5}}(1,2)^\top$, $v_1 = 1$.  So
$A^+ = v_1 \frac{1}{\sigma_1} u_1^\top = \frac{1}{5}(1, 2)
= \frac{A^\top}{A^\top A}$.  For $b = (3, 1)^\top$, the least-squares
solution is $x^* = A^+ b = (3 + 2)/5 = 1$, and
$Ax^* = (1, 2)^\top$ is the projection of $b$ onto $\range(A)$.
\end{example}

\begin{rlconnection}[title={SVD in Dimensionality Reduction and Linear Function Approximation}]
In reinforcement learning with linear function approximation,
the value function is approximated as $\hat{V} = \Phi w$ where
$\Phi \in \R^{|\mathcal{S}| \times d}$ is a feature matrix.
The least-squares solution $w^* = \Phi^+ V^\pi$ uses the
pseudoinverse.  When $d$ is large, the SVD of $\Phi$ reveals which
feature directions carry the most information (large singular values)
and which are redundant (small singular values).  Truncating to the
top $k$ singular values gives the best rank-$k$ feature approximation
by the Eckart--Young theorem, which is precisely what PCA computes
when $\Phi$ contains centered data.
\end{rlconnection}

With the SVD established, we can now define and analyze the matrix
norms that measure the ``size'' of a linear map.

%% ============================================================
%% SECTION 5: TRACE AND MATRIX NORMS
%% ============================================================
\section{Trace and Matrix Norms}
\label{sec:norms}

\subsection{Trace: Definition and Properties}

\begin{definition}[Trace]
\label{def:trace}
The \emph{trace} of a square matrix $A \in \R^{n \times n}$ is
$\tr(A) = \sum_{i=1}^n a_{ii}$.
\end{definition}

\begin{proposition}[Properties of the Trace]
\label{prop:trace_properties}
For matrices $A, B \in \R^{n \times n}$, $C \in \R^{m \times n}$,
$D \in \R^{n \times m}$, and $\alpha \in \R$:
\begin{enumerate}[nosep]
  \item \textbf{Linearity:} $\tr(\alpha A + B) = \alpha \tr(A) + \tr(B)$.
  \item \textbf{Cyclic property:} $\tr(CD) = \tr(DC)$.
  \item \textbf{Similarity invariance:} If $P$ is invertible, then
        $\tr(P^{-1}AP) = \tr(A)$.
  \item \textbf{Eigenvalue sum:} $\tr(A) = \sum_{i=1}^n \lambda_i$
        where $\lambda_i$ are the eigenvalues of $A$ (counted with
        algebraic multiplicity).
  \item \textbf{Frobenius inner product:}
        $\tr(A^\top B) = \sum_{i,j} a_{ij} b_{ij}$.
\end{enumerate}
\end{proposition}

\begin{proof}
(1) Direct from the definition:
$\tr(\alpha A + B) = \sum_i (\alpha a_{ii} + b_{ii})
= \alpha \sum_i a_{ii} + \sum_i b_{ii}$.

(2) Let $C = (c_{ik}) \in \R^{m \times n}$, $D = (d_{kj}) \in \R^{n \times m}$.
Then $\tr(CD) = \sum_{i=1}^m \sum_{k=1}^n c_{ik} d_{ki}
= \sum_{k=1}^n \sum_{i=1}^m d_{ki} c_{ik} = \tr(DC)$.

(3) Follows from (2): $\tr(P^{-1}AP) = \tr(APP^{-1}) = \tr(A)$.

(4) Over $\C$, every matrix is similar to an upper triangular matrix $T$
(Schur decomposition; see Lesson~6).  Then $\tr(A) = \tr(T) = \sum_i t_{ii}
= \sum_i \lambda_i$.

(5) $\tr(A^\top B) = \sum_i (A^\top B)_{ii} = \sum_i \sum_j a_{ji} b_{ji}
= \sum_{i,j} a_{ij} b_{ij}$.
\end{proof}

\begin{example}[Trace as an Inner Product on Matrices]
\label{ex:trace_inner_product}
The \emph{Frobenius inner product} on $\R^{m \times n}$ is
$\langle A, B \rangle_F = \tr(A^\top B)$.  It satisfies all inner
product axioms: symmetry ($\tr(A^\top B) = \tr(B^\top A)$),
bilinearity (from linearity of trace), and positive definiteness
($\tr(A^\top A) = \sum_{ij} a_{ij}^2 \geq 0$, with equality iff
$A = 0$).  This turns $\R^{m \times n}$ into an inner product space
of dimension $mn$.
\end{example}

\subsection{Operator Norm}

\begin{definition}[Operator Norm]
\label{def:op_norm}
For $A \in \R^{m \times n}$, the \emph{operator norm} (or
\emph{spectral norm}, or \emph{induced 2-norm}) is
\[
  \|A\|_{\mathrm{op}} = \max_{\|x\| = 1} \|Ax\|
  = \max_{x \neq 0} \frac{\|Ax\|}{\|x\|},
\]
where $\|\cdot\|$ denotes the Euclidean norm on $\R^n$ and $\R^m$.
\end{definition}

\begin{proposition}[Operator Norm Equals $\sigma_1$]
\label{prop:op_norm_sigma}
$\|A\|_{\mathrm{op}} = \sigma_1(A)$, the largest singular value of $A$.
\end{proposition}

\begin{proof}
For any unit vector $x$, write $x = \sum_j c_j v_j$ with
$\sum_j c_j^2 = 1$ (in the right singular vector basis).  Then
$Ax = \sum_{j=1}^r c_j \sigma_j u_j$ and
$\|Ax\|^2 = \sum_{j=1}^r c_j^2 \sigma_j^2 \leq \sigma_1^2 \sum_j c_j^2
= \sigma_1^2$.
Equality holds at $x = v_1$.
\end{proof}

\subsection{Frobenius Norm}

\begin{definition}[Frobenius Norm]
\label{def:frobenius_norm}
The \emph{Frobenius norm} of $A \in \R^{m \times n}$ is
\[
  \|A\|_F = \sqrt{\sum_{i=1}^m \sum_{j=1}^n a_{ij}^2}
  = \sqrt{\tr(A^\top A)}.
\]
\end{definition}

\begin{proposition}[Frobenius Norm from Singular Values]
\label{prop:frob_singular}
$\|A\|_F = \sqrt{\sigma_1^2 + \cdots + \sigma_r^2}$ where
$\sigma_1, \ldots, \sigma_r$ are the nonzero singular values of $A$.
\end{proposition}

\begin{proof}
$\|A\|_F^2 = \tr(A^\top A) = \tr(V\Sigma^\top\Sigma V^\top)
= \tr(\Sigma^\top\Sigma) = \sum_i \sigma_i^2$,
using the cyclic property of trace and the fact that
$\Sigma^\top \Sigma = \diag(\sigma_1^2, \ldots, \sigma_n^2)$.
\end{proof}

\begin{example}[Computing Both Norms]
\label{ex:norms_computation}
Let $A = \begin{pmatrix} 1 & 2 \\ 0 & 1 \end{pmatrix}$.
The Frobenius norm is $\|A\|_F = \sqrt{1 + 4 + 0 + 1} = \sqrt{6}$.
For the operator norm, $A^\top A = \begin{pmatrix} 1 & 2 \\ 2 & 5 \end{pmatrix}$
with eigenvalues $\mu = 3 \pm \sqrt{5}$.  So
$\sigma_1 = \sqrt{3 + \sqrt{5}} \approx 2.29$ and
$\|A\|_{\mathrm{op}} = \sigma_1$.  We verify:
$\sigma_1^2 + \sigma_2^2 = (3 + \sqrt{5}) + (3 - \sqrt{5}) = 6 = \|A\|_F^2$.
\end{example}

\subsection{Norm Relationships}

\begin{theorem}[Norm Inequalities]
\label{thm:norm_ineq}
For $A \in \R^{m \times n}$ with $\rank(A) = r$:
\[
  \|A\|_{\mathrm{op}} \leq \|A\|_F \leq \sqrt{r}\,\|A\|_{\mathrm{op}}.
\]
\end{theorem}

\begin{proof}
For the left inequality:
$\|A\|_{\mathrm{op}}^2 = \sigma_1^2 \leq \sigma_1^2 + \cdots + \sigma_r^2
= \|A\|_F^2$.

For the right inequality:
$\|A\|_F^2 = \sum_{i=1}^r \sigma_i^2 \leq r \cdot \sigma_1^2
= r \cdot \|A\|_{\mathrm{op}}^2$.

Both bounds are tight: the left is tight when $\rank(A) = 1$
(all but one singular value is zero), and the right is tight when
$\sigma_1 = \cdots = \sigma_r$ (e.g., $A = I_r$ padded with zeros).
\end{proof}

\subsection{Submultiplicativity}

\begin{theorem}[Submultiplicativity]
\label{thm:submult}
For $A \in \R^{m \times p}$ and $B \in \R^{p \times n}$:
\begin{enumerate}[nosep]
  \item $\|AB\|_{\mathrm{op}} \leq \|A\|_{\mathrm{op}} \cdot \|B\|_{\mathrm{op}}$.
  \item $\|AB\|_F \leq \|A\|_{\mathrm{op}} \cdot \|B\|_F$.
  \item $\|AB\|_F \leq \|A\|_F \cdot \|B\|_F$.
\end{enumerate}
\end{theorem}

\begin{proof}
(1) For any unit vector $x$:
$\|ABx\| \leq \|A\|_{\mathrm{op}} \|Bx\|
\leq \|A\|_{\mathrm{op}} \|B\|_{\mathrm{op}} \|x\|
= \|A\|_{\mathrm{op}} \|B\|_{\mathrm{op}}$.
Taking the maximum over $\|x\| = 1$ gives the result.

(2) Let $b_1, \ldots, b_n$ be the columns of $B$.  Then
$\|AB\|_F^2 = \sum_{j=1}^n \|Ab_j\|^2
\leq \|A\|_{\mathrm{op}}^2 \sum_{j=1}^n \|b_j\|^2
= \|A\|_{\mathrm{op}}^2 \|B\|_F^2$.

(3) By the Cauchy--Schwarz inequality on the Frobenius inner product:
$(AB)_{ij} = \sum_k a_{ik} b_{kj}$, so
$\|AB\|_F^2 = \sum_{i,j} \bigl(\sum_k a_{ik} b_{kj}\bigr)^2$.
Alternatively, from (2): $\|AB\|_F \leq \|A\|_{\mathrm{op}} \|B\|_F
\leq \|A\|_F \|B\|_F$ using
$\|A\|_{\mathrm{op}} \leq \|A\|_F$ (Theorem~\ref{thm:norm_ineq}).
\end{proof}

\subsection{Condition Number}

\begin{definition}[Condition Number]
\label{def:condition_number}
For an invertible matrix $A \in \R^{n \times n}$, the
\emph{condition number} (with respect to the operator norm) is
\[
  \kappa(A) = \|A\|_{\mathrm{op}} \cdot \|A^{-1}\|_{\mathrm{op}}
  = \frac{\sigma_1(A)}{\sigma_n(A)}.
\]
\end{definition}

\begin{theorem}[Sensitivity of Linear Systems]
\label{thm:sensitivity}
Let $A \in \R^{n \times n}$ be invertible and let $Ax = b$.  If
$b$ is perturbed to $b + \delta b$ with solution $x + \delta x$, then
\[
  \frac{\|\delta x\|}{\|x\|}
  \leq \kappa(A)\, \frac{\|\delta b\|}{\|b\|}.
\]
\end{theorem}

\begin{proof}
From $A(x + \delta x) = b + \delta b$, we get $A\,\delta x = \delta b$,
so $\delta x = A^{-1}\delta b$.  Hence
$\|\delta x\| \leq \|A^{-1}\|_{\mathrm{op}} \|\delta b\|$.
Also $\|b\| = \|Ax\| \leq \|A\|_{\mathrm{op}} \|x\|$, so
$1/\|x\| \leq \|A\|_{\mathrm{op}}/\|b\|$.  Combining:
\[
  \frac{\|\delta x\|}{\|x\|}
  \leq \|A^{-1}\|_{\mathrm{op}} \|A\|_{\mathrm{op}}
       \frac{\|\delta b\|}{\|b\|}
  = \kappa(A) \frac{\|\delta b\|}{\|b\|}. \qedhere
\]
\end{proof}

\begin{example}[Condition Number and Sensitivity]
\label{ex:condition}
The matrix $A = \begin{pmatrix} 1 & 0 \\ 0 & 0.001 \end{pmatrix}$
has $\sigma_1 = 1$, $\sigma_2 = 0.001$, so $\kappa(A) = 1000$.
A $0.1\%$ perturbation in $b$ can cause up to a $100\%$ change in $x$.
In contrast, $A = I$ has $\kappa = 1$ (perfectly conditioned).
\end{example}

\begin{warning}[title={Condition Number and Numerical Stability}]
The condition number $\kappa(A)$ is an intrinsic property of the
problem $Ax = b$, not of any particular algorithm.  A large condition
number means the problem is inherently sensitive to perturbations.
In machine learning, the condition number of the Hessian
$\kappa(H_f)$ determines the convergence rate of gradient descent:
the number of iterations scales as $O(\kappa \log(1/\varepsilon))$.
This motivates preconditioning techniques and adaptive methods that
implicitly rescale the problem to reduce $\kappa$.
\end{warning}

With norms providing quantitative measures of matrix size and
sensitivity, we now develop the matrix exponential, which connects
linear algebra to continuous-time dynamics.

%% ============================================================
%% SECTION 6: MATRIX EXPONENTIAL
%% ============================================================
\section{Matrix Exponential}
\label{sec:matrix_exp}

\subsection{Definition and Convergence}

\begin{definition}[Matrix Exponential]
\label{def:matrix_exp}
For $A \in \R^{n \times n}$, the \emph{matrix exponential} is
\[
  e^A = \exp(A) = \sum_{k=0}^{\infty} \frac{A^k}{k!}
  = I + A + \frac{A^2}{2!} + \frac{A^3}{3!} + \cdots.
\]
\end{definition}

\begin{proposition}[Convergence of the Matrix Exponential]
\label{prop:exp_convergence}
The series $\sum_{k=0}^\infty A^k / k!$ converges absolutely (in
any matrix norm) for every $A \in \R^{n \times n}$.
\end{proposition}

\begin{proof}
Using submultiplicativity (Theorem~\ref{thm:submult}),
$\|A^k/k!\|_{\mathrm{op}} \leq \|A\|_{\mathrm{op}}^k / k!$.
The scalar series $\sum_k \|A\|_{\mathrm{op}}^k / k!
= e^{\|A\|_{\mathrm{op}}} < \infty$ converges, so the matrix
series converges absolutely.  (Absolute convergence in a complete normed
space implies convergence.)
\end{proof}

\subsection{Properties}

\begin{theorem}[Properties of the Matrix Exponential]
\label{thm:exp_properties}
Let $A, B \in \R^{n \times n}$.
\begin{enumerate}[nosep]
  \item $e^{0} = I$.
  \item If $AB = BA$, then $e^{A+B} = e^A e^B$.
  \item $e^A$ is always invertible, with $(e^A)^{-1} = e^{-A}$.
  \item $\frac{d}{dt} e^{tA} = A e^{tA} = e^{tA} A$.
  \item If $A = PDP^{-1}$ (diagonalizable), then
        $e^A = P\, \diag(e^{\lambda_1}, \ldots, e^{\lambda_n})\, P^{-1}$.
\end{enumerate}
\end{theorem}

\begin{proof}
(1) $e^0 = \sum_k 0^k/k! = I$ (where $0^0 = I$ by convention).

(2) When $AB = BA$, the binomial theorem applies:
$(A + B)^k = \sum_{j=0}^k \binom{k}{j} A^j B^{k-j}$.  Then
\begin{align*}
  e^{A+B} &= \sum_{k=0}^\infty \frac{(A+B)^k}{k!}
  = \sum_{k=0}^\infty \sum_{j=0}^k \frac{A^j}{j!} \frac{B^{k-j}}{(k-j)!}
  = \Bigl(\sum_{j=0}^\infty \frac{A^j}{j!}\Bigr)
    \Bigl(\sum_{\ell=0}^\infty \frac{B^\ell}{\ell!}\Bigr)
  = e^A e^B,
\end{align*}
where the rearrangement is justified by absolute convergence.

(3) Since $A$ commutes with $-A$, part (2) gives
$e^A e^{-A} = e^{A-A} = e^0 = I$.

(4) Differentiating term by term (justified by uniform convergence on
compact intervals):
\[
  \frac{d}{dt} e^{tA} = \frac{d}{dt} \sum_{k=0}^\infty \frac{t^k A^k}{k!}
  = \sum_{k=1}^\infty \frac{t^{k-1} A^k}{(k-1)!}
  = A \sum_{k=1}^\infty \frac{t^{k-1} A^{k-1}}{(k-1)!}
  = A e^{tA}.
\]
The identity $e^{tA} A$ follows since $A$ commutes with $e^{tA}$
(each term $t^k A^k / k!$ commutes with $A$).

(5) $A^k = PD^kP^{-1}$, so
$e^A = P\bigl(\sum_k D^k/k!\bigr)P^{-1}
= P\,\diag(e^{\lambda_1}, \ldots, e^{\lambda_n})\,P^{-1}$.
\end{proof}

\begin{example}[Matrix Exponential of a Diagonal Matrix]
\label{ex:exp_diagonal}
If $A = \diag(-1, 2)$, then
$e^{tA} = \diag(e^{-t}, e^{2t})$.  As $t \to \infty$, the first
component decays to $0$ and the second grows without bound.  The
eigenvalues $\{-1, 2\}$ determine the long-time behavior.
\end{example}

\begin{example}[Matrix Exponential and Rotation]
\label{ex:exp_rotation}
Let $A = \begin{pmatrix} 0 & -\theta \\ \theta & 0 \end{pmatrix}$
for $\theta \in \R$.  Since $A$ is skew-symmetric ($A^\top = -A$),
$e^A$ is orthogonal.  Using $A^2 = -\theta^2 I$, the series groups as
\[
  e^A = I \cos\theta + \frac{A}{\theta}\sin\theta
  = \begin{pmatrix} \cos\theta & -\sin\theta \\ \sin\theta & \cos\theta \end{pmatrix},
\]
which is a rotation by angle $\theta$.
\end{example}

\begin{example}[Linear ODE Solution]
\label{ex:ode_solution}
The system $\dot{x}(t) = Ax(t)$ with initial condition $x(0) = x_0$
has the unique solution $x(t) = e^{tA} x_0$.  This follows from
property (4): $\frac{d}{dt}(e^{tA}x_0) = Ae^{tA}x_0 = Ax(t)$.
In continuous-time Markov chains, the transition rate matrix $Q$
(with $q_{ij} \geq 0$ for $i \neq j$ and row sums zero) gives the
transition probability matrix $P(t) = e^{tQ}$.
\end{example}

We now develop the final topic: the systematic computation of
derivatives of matrix-valued expressions.

%% ============================================================
%% SECTION 7: MATRIX CALCULUS
%% ============================================================
\section{Matrix Calculus}
\label{sec:matrix_calculus}

\subsection{Gradient, Jacobian, and Hessian}

\begin{definition}[Gradient]
\label{def:gradient}
Let $f\colon \R^n \to \R$ be differentiable at $x$.  The
\emph{gradient} of $f$ at $x$ is the unique vector
$\nabla f(x) \in \R^n$ satisfying
\[
  f(x + h) = f(x) + \nabla f(x)^\top h + o(\|h\|)
  \quad \text{as } \|h\| \to 0.
\]
In coordinates, $[\nabla f(x)]_i = \partial f / \partial x_i$.
\end{definition}

\begin{definition}[Jacobian]
\label{def:jacobian}
Let $g\colon \R^n \to \R^m$ be differentiable at $x$.  The
\emph{Jacobian} of $g$ at $x$ is the matrix $J_g(x) \in \R^{m \times n}$
with $[J_g(x)]_{ij} = \partial g_i / \partial x_j$, satisfying
\[
  g(x + h) = g(x) + J_g(x)\, h + o(\|h\|).
\]
When $m = 1$, $J_g(x) = \nabla g(x)^\top$.
\end{definition}

\begin{definition}[Hessian]
\label{def:hessian}
Let $f\colon \R^n \to \R$ be twice differentiable at $x$.  The
\emph{Hessian} of $f$ at $x$ is $H_f(x) \in \R^{n \times n}$ with
$[H_f(x)]_{ij} = \frac{\partial^2 f}{\partial x_i \partial x_j}$.
If the second partial derivatives are continuous, $H_f(x)$ is symmetric.
The Hessian is the Jacobian of the gradient: $H_f(x) = J_{\nabla f}(x)$.
\end{definition}

\subsection{Fundamental Vector Derivative Formulas}

\begin{theorem}[Basic Gradient Formulas]
\label{thm:basic_gradients}
Let $a \in \R^n$, $A \in \R^{n \times n}$, and $x \in \R^n$.
\begin{enumerate}[nosep]
  \item $\nabla_x(a^\top x) = a$.
  \item $\nabla_x(x^\top A x) = (A + A^\top) x$.
        If $A$ is symmetric, this simplifies to $2Ax$.
  \item $\nabla_x(\|x\|^2) = 2x$.
  \item $\nabla_x(\|Ax - b\|^2) = 2A^\top(Ax - b)$ for $b \in \R^m$,
        $A \in \R^{m \times n}$.
\end{enumerate}
\end{theorem}

\begin{proof}
(1) $f(x + h) = a^\top(x + h) = a^\top x + a^\top h = f(x) + a^\top h$.
Comparing with $f(x) + \nabla f(x)^\top h + o(\|h\|)$ gives
$\nabla f(x) = a$.

(2) $f(x + h) = (x + h)^\top A (x + h)
= x^\top Ax + x^\top Ah + h^\top Ax + h^\top Ah$.
The linear term in $h$ is $x^\top Ah + h^\top Ax = (A^\top x)^\top h + (Ax)^\top h
= (A^\top x + Ax)^\top h = [(A + A^\top)x]^\top h$.  The term $h^\top Ah$
is $o(\|h\|)$.  Hence $\nabla f(x) = (A + A^\top)x$.

(3) $\|x\|^2 = x^\top I x$.  Apply (2) with $A = I$:
$\nabla_x(x^\top x) = 2Ix = 2x$.

(4) Set $g(x) = \|Ax - b\|^2 = (Ax - b)^\top(Ax - b)
= x^\top A^\top A x - 2b^\top Ax + b^\top b$.  By (1) and (2):
$\nabla g(x) = (A^\top A + A^\top A)x - 2A^\top b
= 2A^\top A x - 2A^\top b = 2A^\top(Ax - b)$.
\end{proof}

\begin{example}[Gradient of a Quadratic Loss]
\label{ex:quadratic_loss}
The least-squares loss $L(w) = \frac{1}{2}\|Xw - y\|^2$ with
$X \in \R^{m \times n}$ and $y \in \R^m$ has gradient
$\nabla L(w) = X^\top(Xw - y)$
(by Theorem~\ref{thm:basic_gradients}(4) with a factor of $1/2$).
Setting $\nabla L = 0$ gives the normal equations
$X^\top X w = X^\top y$, whose minimum-norm solution is $w^* = X^+ y$.
\end{example}

\begin{example}[Gradient of a Regularized Loss]
\label{ex:ridge}
Ridge regression adds an $\ell_2$ penalty:
$L(w) = \frac{1}{2}\|Xw - y\|^2 + \frac{\lambda}{2}\|w\|^2$.
Using Theorem~\ref{thm:basic_gradients}:
$\nabla L(w) = X^\top(Xw - y) + \lambda w$.
Setting to zero: $(X^\top X + \lambda I)w = X^\top y$,
so $w^* = (X^\top X + \lambda I)^{-1} X^\top y$.  The regularization
term $\lambda I$ improves the condition number of $X^\top X$.
\end{example}

\subsection{Trace Derivative Formulas}

When the variable is a matrix $X \in \R^{m \times n}$, the gradient
$\nabla_X f$ is the matrix of partial derivatives
$[\nabla_X f]_{ij} = \partial f / \partial x_{ij}$, characterized by
\[
  f(X + H) = f(X) + \tr(\nabla_X f(X)^\top H) + o(\|H\|_F).
\]

\begin{theorem}[Trace Derivative Formulas]
\label{thm:trace_derivs}
Let $A \in \R^{n \times m}$, $B \in \R^{m \times n}$,
$X \in \R^{m \times n}$.
\begin{enumerate}[nosep]
  \item $\frac{\partial}{\partial X} \tr(AX) = A^\top$.
  \item $\frac{\partial}{\partial X} \tr(X^\top A) = A$.
  \item $\frac{\partial}{\partial X} \tr(X^\top A X)
        = (A + A^\top) X$ \quad (for $X \in \R^{n \times p}$,
        $A \in \R^{n \times n}$).
  \item $\frac{\partial}{\partial X} \tr(AXB) = A^\top B^\top$ \quad
        (for $A \in \R^{p \times m}$, $B \in \R^{n \times p}$).
\end{enumerate}
\end{theorem}

\begin{proof}
(1) $\tr(A(X+H)) = \tr(AX) + \tr(AH)$.  Since
$\tr(AH) = \tr(A^\top{}^\top H) = \sum_{ij} [A^\top]_{ij} h_{ij}
= \tr((A^\top)^\top H)$, the gradient is $A^\top$.

(2) $\tr((X+H)^\top A) = \tr(X^\top A) + \tr(H^\top A)$.
Now $\tr(H^\top A) = \sum_{ij} h_{ij} a_{ij} = \tr(A^\top H)$, so
the gradient is $A$.

(3) $\tr((X+H)^\top A(X+H)) = \tr(X^\top AX) + \tr(H^\top AX)
+ \tr(X^\top AH) + \tr(H^\top AH)$.  The linear part is
$\tr(H^\top AX) + \tr(X^\top AH)
= \tr(H^\top AX) + \tr(H^\top A^\top X)
= \tr(H^\top (A + A^\top)X)$, giving gradient $(A + A^\top)X$.

(4) $\tr(A(X+H)B) = \tr(AXB) + \tr(AHB)$.  Using the cyclic property:
$\tr(AHB) = \tr(BAHB \cdot B^{-1}) = \tr(BAH)$ -- more directly,
$\tr(AHB) = \tr(B A H) = \tr((A^\top B^\top)^\top H)$, so the gradient
is $A^\top B^\top$.
\end{proof}

\begin{example}[Gradient of $\|AX - B\|_F^2$]
\label{ex:matrix_ls}
Let $f(X) = \|AX - B\|_F^2 = \tr((AX - B)^\top(AX - B))$.  Expanding:
$f(X) = \tr(X^\top A^\top A X) - 2\tr(B^\top AX) + \tr(B^\top B)$.
Using Theorem~\ref{thm:trace_derivs}:
$\nabla_X f = 2A^\top A X - 2A^\top B = 2A^\top(AX - B)$.
This extends the vector least-squares gradient to the matrix case.
\end{example}

\subsection{The Chain Rule}

\begin{theorem}[Chain Rule for Compositions]
\label{thm:chain_rule}
Let $g\colon \R^n \to \R^m$ be differentiable at $x$ and
$f\colon \R^m \to \R$ be differentiable at $g(x)$.  Then the
composition $h = f \circ g\colon \R^n \to \R$ is differentiable at $x$
with
\[
  \nabla h(x) = J_g(x)^\top \nabla f(g(x)).
\]
More generally, if $f\colon \R^m \to \R^p$, then $J_{f \circ g}(x) = J_f(g(x))\, J_g(x)$.
\end{theorem}

\begin{proof}
By differentiability of $g$ and $f$:
\begin{align*}
  h(x + \delta) &= f(g(x + \delta))
  = f(g(x) + J_g(x)\delta + o(\|\delta\|)) \\
  &= f(g(x)) + \nabla f(g(x))^\top [J_g(x)\delta + o(\|\delta\|)]
     + o(\|J_g(x)\delta + o(\|\delta\|)\|) \\
  &= f(g(x)) + \nabla f(g(x))^\top J_g(x)\delta + o(\|\delta\|) \\
  &= h(x) + [J_g(x)^\top \nabla f(g(x))]^\top \delta + o(\|\delta\|).
\end{align*}
Hence $\nabla h(x) = J_g(x)^\top \nabla f(g(x))$.  The general case
follows by applying this to each component of $f$.
\end{proof}

\begin{keyresult}[title={Matrix Calculus Toolbox for Gradient Computation}]
The results of this section provide the complete toolbox for computing
gradients in machine learning:
\begin{enumerate}[nosep]
  \item \textbf{Vector derivatives:} $\nabla(a^\top x) = a$,
        $\nabla(x^\top Ax) = (A + A^\top)x$ (Theorem~\ref{thm:basic_gradients}).
  \item \textbf{Trace derivatives:} $\partial \tr(AX)/\partial X = A^\top$,
        $\partial \tr(X^\top AX)/\partial X = (A + A^\top)X$
        (Theorem~\ref{thm:trace_derivs}).
  \item \textbf{Chain rule:} $\nabla(f \circ g)(x)
        = J_g(x)^\top \nabla f(g(x))$ (Theorem~\ref{thm:chain_rule}).
\end{enumerate}
Every gradient computation in backpropagation is an iterated application
of these three rules.  The chain rule composes the Jacobians of
individual layers, while the vector and trace formulas compute the
Jacobian of each layer.
\end{keyresult}

\begin{example}[Softmax Cross-Entropy Gradient]
\label{ex:softmax}
In a classification network, the final layer computes logits
$z = Wx + b \in \R^K$, then $p = \mathrm{softmax}(z)$ with
$p_k = e^{z_k}/\sum_j e^{z_j}$, and the loss is
$L = -\log p_y$ (cross-entropy for true class $y$).  By the chain
rule: $\nabla_z L = p - e_y$ (where $e_y$ is the $y$-th standard
basis vector), and then $\nabla_W L = (p - e_y)x^\top$ (a rank-1
outer product).  Each step uses the chain rule
(Theorem~\ref{thm:chain_rule}) applied to the composition of the
linear layer, softmax, and log-loss.
\end{example}

%% ============================================================
%% SECTION 8: EXERCISES
%% ============================================================
\section{Exercises}
\label{sec:exercises}

\begin{exercise}[Quadratic Form Classification]
\label{exer:qf}
\textnormal{(\(*\))} \\
For each matrix, classify the associated quadratic form as positive
definite, positive semidefinite, negative definite, negative
semidefinite, or indefinite:
\begin{enumerate}[nosep,label=(\alph*)]
  \item $A = \begin{pmatrix} 5 & 2 \\ 2 & 1 \end{pmatrix}$.
  \item $B = \begin{pmatrix} -3 & 0 \\ 0 & -1 \end{pmatrix}$.
  \item $C = \begin{pmatrix} 1 & 0 & 0 \\ 0 & 0 & 0 \\ 0 & 0 & -1 \end{pmatrix}$.
\end{enumerate}
\end{exercise}

\begin{exercise}[SVD Computation]
\label{exer:svd}
\textnormal{(\(*\))} \\
Compute the full SVD of $A = \begin{pmatrix} 1 & 1 \\ 1 & -1 \\ 0 & 0 \end{pmatrix}$.
\begin{enumerate}[nosep,label=(\alph*)]
  \item Find $A^\top A$ and its eigenvalues and eigenvectors.
  \item Determine the singular values and right singular vectors.
  \item Construct the left singular vectors and verify $A = U\Sigma V^\top$.
  \item Find the pseudoinverse $A^+$ and verify $A A^+ A = A$.
\end{enumerate}
\end{exercise}

\begin{exercise}[Low-Rank Approximation in RL]
\label{exer:eckart_young}
\textnormal{(\(**\))} \\
In a reinforcement learning problem with $|\mathcal{S}| = 100$ states
and $d = 10$ features, the feature matrix $\Phi \in \R^{100 \times 10}$
has singular values
$\sigma_1 = 5, \sigma_2 = 4, \sigma_3 = 3, \sigma_4 = 2, \sigma_5 = 1,
\sigma_6 = \cdots = \sigma_{10} = 0.01$.
\begin{enumerate}[nosep,label=(\alph*)]
  \item Compute $\kappa(\Phi^\top \Phi)$ and interpret its meaning
        for the least-squares problem $\Phi w \approx V^\pi$.
  \item If we truncate $\Phi$ to rank $k = 5$, what is the Frobenius
        norm of the approximation error $\|\Phi - \Phi_5\|_F$?
  \item What fraction of $\|\Phi\|_F^2$ is captured by the rank-5
        approximation?
\end{enumerate}
\end{exercise}

\begin{exercise}[Matrix Exponential Computation]
\label{exer:exp}
\textnormal{(\(**\))} \\
Let $A = \begin{pmatrix} 1 & 1 \\ 0 & 2 \end{pmatrix}$.
\begin{enumerate}[nosep,label=(\alph*)]
  \item Find $P$ and $D$ such that $A = PDP^{-1}$.
  \item Compute $e^{tA}$ for all $t \in \R$.
  \item Verify that $\frac{d}{dt}e^{tA}\big|_{t=0} = A$.
  \item Solve the ODE $\dot{x} = Ax$ with $x(0) = (1, 0)^\top$.
\end{enumerate}
\end{exercise}

\begin{exercise}[Trace Derivative Derivation]
\label{exer:trace_deriv}
\textnormal{(\(**\))} \\
\begin{enumerate}[nosep,label=(\alph*)]
  \item Prove that $\frac{\partial}{\partial X} \tr(AX^\top BX)
        = A^\top X B^\top + BXA$ for $X \in \R^{m \times n}$,
        $A \in \R^{m \times m}$, $B \in \R^{n \times n}$.
  \item Using part (a) and the identity $\|AX - B\|_F^2 = \tr((AX-B)^\top(AX-B))$,
        re-derive $\nabla_X \|AX - B\|_F^2 = 2A^\top(AX - B)$.
\end{enumerate}
\textit{Hint:} For (a), expand $\tr(A(X+H)^\top B(X+H))$ and isolate
the terms linear in $H$.
\end{exercise}

\begin{exercise}[Gradient of the Log-Determinant]
\label{exer:logdet}
\textnormal{(\(***\))} \\
Let $f(X) = \log\det(X)$ for symmetric positive definite
$X \in \R^{n \times n}$.
\begin{enumerate}[nosep,label=(\alph*)]
  \item Show that $\det(X + tH) = \det(X)\det(I + tX^{-1}H)$
        for small $t$.
  \item Using $\det(I + tM) = 1 + t\,\tr(M) + O(t^2)$, prove that
        $\frac{d}{dt}\big|_{t=0} \log\det(X + tH) = \tr(X^{-1}H)$.
  \item Conclude that $\nabla_X \log\det(X) = X^{-1}$ (in the sense
        of the Frobenius inner product: $\tr((\nabla f)^\top H) = \tr(X^{-1}H)$).
  \item Discuss the relevance of this result to Gaussian distributions,
        where the log-likelihood involves $\log\det(\Sigma)$ for the
        covariance matrix $\Sigma$.
\end{enumerate}
\textit{Hint:} For (b), recall that $\log(1 + \varepsilon) = \varepsilon + O(\varepsilon^2)$.
\end{exercise}

%% ============================================================
%% SECTION 9: SUMMARY
%% ============================================================
\section{Summary}
\label{sec:summary}

\begin{itemize}
  \item A \emph{quadratic form} $Q_A(x) = x^\top Ax$ is classified
        by the eigenvalues of the symmetric matrix $A$: positive definite
        iff all $\lambda_i > 0$, positive semidefinite iff all
        $\lambda_i \geq 0$, indefinite iff $A$ has eigenvalues of
        both signs.  At a critical point, the Hessian quadratic form
        determines the local nature of the extremum.

  \item The \emph{singular value decomposition} $A = U\Sigma V^\top$
        factors any $m \times n$ matrix into orthogonal rotations and
        a diagonal scaling.  The singular values are the square roots
        of the eigenvalues of $A^\top A$, and $\rank(A)$ equals the
        number of nonzero singular values.

  \item The \emph{Eckart--Young theorem} guarantees that truncating the
        SVD at rank $k$ gives the best rank-$k$ approximation in both
        operator and Frobenius norm.  This is the theoretical
        foundation of PCA.

  \item The \emph{Moore--Penrose pseudoinverse} $A^+$ is the unique
        matrix satisfying the four Penrose conditions.  The vector
        $A^+ b$ is the minimum-norm least-squares solution to $Ax = b$.

  \item The \emph{operator norm} $\|A\|_{\mathrm{op}} = \sigma_1$ and
        the \emph{Frobenius norm} $\|A\|_F = (\sum \sigma_i^2)^{1/2}$
        satisfy $\|A\|_{\mathrm{op}} \leq \|A\|_F \leq \sqrt{r}\,\|A\|_{\mathrm{op}}$.
        Both are submultiplicative.

  \item The \emph{condition number} $\kappa(A) = \sigma_1/\sigma_n$
        bounds the relative amplification of perturbations in linear
        systems.  Large $\kappa$ means the system is ill-conditioned.

  \item The \emph{matrix exponential} $e^A = \sum_{k=0}^\infty A^k/k!$
        converges for every square matrix.  It satisfies
        $\frac{d}{dt}e^{tA} = Ae^{tA}$, solving the linear ODE
        $\dot{x} = Ax$ with $x(t) = e^{tA}x_0$.

  \item \emph{Matrix calculus} provides systematic rules:
        $\nabla(a^\top x) = a$,
        $\nabla(x^\top Ax) = (A + A^\top)x$,
        $\partial\tr(AX)/\partial X = A^\top$,
        and the chain rule $\nabla(f \circ g) = J_g^\top \nabla f$.

  \item The trace function defines an inner product on matrices
        ($\langle A, B \rangle_F = \tr(A^\top B)$) and provides a
        unified framework for expressing matrix derivatives.
\end{itemize}

\paragraph{Capstone summary: Linear Algebra subdomain (Lessons 5--8).}

\begin{table}[h]
\centering
\begin{tabular}{@{} c l p{7.2cm} @{}}
\toprule
\textbf{Lesson} & \textbf{Title} & \textbf{Key Results} \\
\midrule
5 & Vector Spaces and Linear Maps &
    Vector space axioms, subspaces, bases and dimension,
    linear maps, rank--nullity theorem, matrix representation,
    change of basis \\
6 & Determinants and Eigenvalues &
    Determinant (cofactor, Leibniz), multiplicativity,
    eigenvalues, characteristic polynomial, algebraic and
    geometric multiplicity, diagonalization, Cayley--Hamilton \\
7 & Inner Products and Spectral Theory &
    Inner products, Cauchy--Schwarz, Gram--Schmidt,
    orthogonal projection and best approximation, spectral
    theorem (real symmetric), positive definiteness, Cholesky,
    Rayleigh quotient, Courant--Fischer \\
8 & SVD, Norms, and Matrix Calculus &
    Quadratic forms, SVD theorem, Eckart--Young, pseudoinverse,
    operator and Frobenius norms, condition number, matrix
    exponential, gradient/Jacobian/Hessian, trace derivatives,
    chain rule \\
\bottomrule
\end{tabular}
\caption{Linear Algebra subdomain arc: Lessons 5--8.}
\label{tab:la_summary}
\end{table}

%% ============================================================
%% SECTION 10: FURTHER READING
%% ============================================================
\section{Further Reading}
\label{sec:further_reading}

\begin{itemize}
  \item \textbf{Phase 01, Functional Analysis} (Lessons on metric,
        normed, and Hilbert spaces) generalizes the finite-dimensional
        theory of this lesson to infinite dimensions.  Inner products
        become Hilbert space inner products, the spectral theorem
        extends to compact self-adjoint operators, and the operator norm
        becomes the fundamental tool for analyzing bounded linear maps
        between Banach spaces.

  \item \textbf{Phase 01, Convex Analysis and Optimization} uses
        positive definiteness and matrix calculus extensively.
        The Hessian condition $H_f(x) \succ 0$ characterizes strict
        convexity, and gradient formulas derived here are the starting
        point for convergence analysis of first-order methods.

  \item \citet{magnus2019} is the definitive reference for matrix
        differential calculus, developing the theory via the
        \emph{vec operator} and Kronecker products.  Chapters 8--13
        cover derivatives of eigenvalues, determinants, and matrix
        inverses.

  \item \citet{petersen2012} (the ``Matrix Cookbook'') provides a
        comprehensive catalog of matrix identities, derivatives, and
        decompositions organized for quick reference.  It is an
        indispensable resource for implementing gradient computations
        in neural networks and RL algorithms.

  \item \citet[Chapters 5--7]{horn2013} develops the SVD, matrix
        norms, and perturbation theory (Weyl inequalities, Davis--Kahan
        theorem) in far greater depth, with applications to
        low-rank matrix completion and spectral clustering.

  \item \textbf{Historical note.}  The singular value decomposition was
        independently discovered by Beltrami (1873) and Jordan (1874)
        for bilinear forms, and later placed in modern matrix form by
        Eckart and Young (1936).  The name ``singular values'' was
        coined by Smithies (1937) in the context of integral equations.
        Matrix calculus was systematized by \citet{magnus2019}
        (first edition 1988), building on earlier work by
        Dwyer and Macphail (1948).
\end{itemize}

\bibliography{linear-algebra-refs}
\end{document}
